% Copyright 2026  Ed Bueler

\documentclass[10pt,
               svgnames,
               hyperref={colorlinks,
                         citecolor=DeepPink4,
                         linkcolor=FireBrick,
                         urlcolor=blue},
               usepdftitle=false]{beamer}

\usetheme{Madrid}
\usecolortheme{beaver}
\setbeamercovered{transparent}
\setbeamerfont{frametitle}{size=\large}
\setbeamercolor*{block title}{bg=red!10}
\setbeamercolor*{block body}{bg=red!5}

\usepackage[english]{babel}
\usepackage[latin1]{inputenc}
\usepackage{times}
\usepackage[T1]{fontenc}
% Or whatever. Note that the encoding and the font should match. If T1
% does not look nice, try deleting the line with the fontenc.

\usepackage{fancyvrb,empheq,bm,xspace,dsfont}
%\usepackage{minted}
\usepackage{tikz}

% If you wish to uncover everything in a step-wise fashion, uncomment
% the following command: 
%\beamerdefaultoverlayspecification{<+->}

\newcommand{\bb}{\mathbf{b}}
\newcommand{\bc}{\mathbf{c}}
\newcommand{\bbf}{\mathbf{f}}
\newcommand{\bl}{\bm{\ell}}
\newcommand{\br}{\mathbf{r}}
\newcommand{\bs}{\mathbf{s}}
\newcommand{\bx}{\mathbf{x}}
\newcommand{\by}{\mathbf{y}}
\newcommand{\bv}{\mathbf{v}}
\newcommand{\bu}{\mathbf{u}}
\newcommand{\bw}{\mathbf{w}}

\newcommand{\cB}{\mathcal{B}}
\newcommand{\cM}{\mathcal{M}}
\newcommand{\cR}{\mathcal{R}}

\newcommand{\bzero}{\bm{0}}

\newcommand{\CC}{\mathbb{C}}
\newcommand{\NN}{\mathbb{N}}
\newcommand{\RR}{\mathbb{R}}
\newcommand{\ZZ}{\mathbb{Z}}

\newcommand{\one}{\mathds{1}}

\newcommand{\range}{\operatorname{range}}
\newcommand{\Span}{\operatorname{span}}

\newcommand{\Matlab}{\textsc{Matlab}\xspace}
\newcommand{\eps}{\epsilon}

\newcommand{\ds}{\displaystyle}

\newcommand{\ip}[2]{\left<#1,#2\right>}

\newcommand{\twovect}[4]{\ensuremath{{#1}_{#2} =
                            \begin{bmatrix} #3 \\ #4 \end{bmatrix}}}

\newcommand{\rbullet}{{\color{FireBrick} \bullet}}

\newcommand*\circled[1]{\tikz[baseline=(char.base)]{
            \node[shape=circle,draw,inner sep=2pt] (char) {#1};}}

\newcommand{\textbook}{Saxe, \emph{Beginning Functional Analysis}}

\DefineVerbatimEnvironment{mVerb}{Verbatim}{numbersep=2mm,
frame=lines,framerule=0.1mm,framesep=2mm,xleftmargin=4mm,fontsize=\scriptsize}

\newcommand{\itemnum}[1]{\setcounter{enumi}{#1}\usebeamertemplate{enumerate item}}



\title{Integration by parts and weak derivatives}

\subtitle{a calculation for week 9}

\author{Ed Bueler}

\institute[]{UAF Math 617 Functional Analysis}

\date{Spring 2026}


\begin{document}
\beamertemplatenavigationsymbolsempty

\begin{frame}
  \maketitle
\end{frame}


\begin{frame}{Outline}
  \tableofcontents[hideallsubsections]
\end{frame}

\AtBeginSection[]
{% do not put toc here this time
}


\section{the fundamental theorem of calculus in $\RR^1$}

\begin{frame}{indefinite integral of an $L^1$ function}

\begin{itemize}
\item for $a<b$ real numbers, recall that
	$$L^1([a,b]) = \Big\{f:[a,b]\to\RR\,\big|\,f \text{ is measurable and } \int_a^b |f(x)|\,dx <\infty\Big\}$$

    \begin{itemize}
    \item[$\circ$] remember this is actually a \emph{lie}; the elements of $L^1$ are \emph{equivalence classes} up to almost everywhere
    \end{itemize}
\item what happens when you integrate $f\in L^1$ to get a new function?:
	$$F(x) = \int_a^x f(t)\,dt$$

    \begin{itemize}
    \item[$\circ$] this is the \emph{indefinite integral}
    \end{itemize}
\item does the fundamental theorem of calculus (FTC) apply?:
	$$\frac{d}{dx}\left(\int_a^x f(t)\,dt\right) \stackrel{?}{=} f(x)$$

    \begin{itemize}
    \item[$\circ$] that is, is $F'(x)=f(x)$?
    \end{itemize}
\end{itemize}
\end{frame}


\begin{frame}{fundamental theorem of calculus (FTC)}

\begin{theorem}[easy direction]
If $f\in L^1([a,b])$ then $F(x) = \int_a^x f(t)\,dt$ is continuous, and it is differentiable almost everywhere, and $F'=f$ almost everywhere.
\end{theorem}

\begin{proof}<2>[{\small Proof \alert{when we assume $f$ is continuous}}]
{\footnotesize
If $f\in C([a,b])$ then it is bounded: $|f|\le M$.  Let $\eps>0$.  Suppose $|x-y|<\delta=\frac{\eps}{M}$.  Then
	$$|F(x)-F(y)| = \left|\int_x^y f(t)\,dt\right| \le \int_x^y |f(t)|\,dt \le M |x-y| < M\delta = \eps,$$
so $F$ is continuous.  Since $f$ is continuous at $x$ then there is $\delta>0$ so that $|f(t)-f(x)|<\eps$ if $|t-x|<\delta$.  Choose any $0<h<\delta$.  Then
\begin{align*}
\left|\frac{F(x+h)-F(x)}{h} - f(x)\right| &= \left|\frac{1}{h} \int_x^{x+h} f(t)\,dt -  \frac{1}{h} \int_x^{x+h} f(x)\,dt\right| \\
&\le \frac{1}{h} \int_x^{x+h} |f(t)-f(x)|\,dt < \frac{1}{h} \int_x^{x+h} \eps\,dt = \eps
\end{align*}
(A similar argument handles $-\delta<h<0$.)  This shows $f(x)=\lim_{h\to 0} \frac{F(x+h)-F(x)}{h}$; in particular the two-sided limit exists at every $x$.
}
\end{proof}
\end{frame}


\begin{frame}{picture of the proof}

\noindent picture:

\vspace{70mm}
\end{frame}


\begin{frame}{how the integral is well-behaved when $f\in L^1$}

\begin{lemma}[``absolute continuity of the Lebesgue integral'']
Suppose $f\in L^1([a,b])$ and $\eps>0$.  There exists $\delta>0$ so that if $m(E)<\delta$ then
	$$\int_E |f|\,dm < \eps.$$
\end{lemma}

\begin{proof}[{\small Proof.}]
{\footnotesize
Let $g_k = \min\{|f|, k\}$.  This defines a sequence $0\le g_k \le |f|$ so that $g_k \nearrow |f|$, but where $0 \le g_k \le k$ so $g_k$ is bounded.  By the Monotone Convergence Theorem,
    $$\int_{[a,b]} |f|\,dm = \lim_{k\to\infty} \int_{[a,b]} g_k\,dm$$
Given $\epsilon > 0$, choose $K$ large enough that $\int_{[a,b]} (|f| - g_K)\,dm < \epsilon/2$.  Then for any measurable $E \subset [a,b]$:
    $$\int_E |f|\,dm = \int_E g_K\,dm + \int_E (|f| - g_K) \le \int_E K\,dm + \int_{[a,b]} (|f| - g_K) \le K \cdot m(E) + \frac{\epsilon}{2}$$
Now set $\delta = \eps/(2K)$. If $m(E) < \delta$ then $\int_E |f|\,dm < \eps$.
}
\end{proof}
\end{frame}


\begin{frame}{picture of the absolute continuity of the Lebesgue integral}

\noindent picture:

\vspace{70mm}
\end{frame}


\begin{frame}{classical definition of absolute continuity}

\begin{definition}
a continuous function $f:[a,b]\to\RR$ is \emph{absolutely continuous} if for all $\eps>0$ there is $\delta>0$ so that if we have a finite collection of disjoint intervals $[\alpha_i,\beta_i]$, with total length less than $\delta$:
	$$\sum_{i=1}^n (\beta_i - \alpha_i) <\delta,$$
then the variation of $f$ on those intervals is less than $\eps$:
	$$\sum_{i=1}^n |f(\alpha_i) - f(\beta_i)| < \eps$$
\end{definition}
\end{frame}


\begin{frame}{picture of absolute continuity}

\noindent picture:

\vspace{70mm}
\end{frame}


\begin{frame}{flavors of continuity}

\begin{itemize}
\item uniform continuity $\implies$ continuous at $x\in[a,b]$
\item absolute continuity $\implies$ uniform continuity

\begin{definition}
a continuous function $f:[a,b]\to\RR$ is \emph{Lipschitz continuous} if there is $L\ge 0$ so that
	$$|f(x)-f(y)| \le L|x-y|$$
for all $x,y\in [a,b]$
\end{definition}

\item Lipschitz continuous $\implies$ absolute continuity
\item $C^1([a,b])$ $\implies$ Lipschitz continuous
\end{itemize}
\end{frame}


\begin{frame}{the FTC again}

\begin{theorem}[easy direction?]
If $f\in L^1([a,b])$ then $F(x) = \int_a^x f(t)\,dt$ is \alert{absolutely continuous}, and it is differentiable almost everywhere, and $F'=f$ almost everywhere.
\end{theorem}

\begin{proof}[{\small sketch of the proof}]
{\footnotesize
Approximate $f$ in $L^1$ by a continuous function $g$.  Show that the difference $f-g$, which can be large, can't be large on a big set.  (Vitali covering lemma.)  Conclude.
}
\end{proof}

\begin{theorem}[other direction]
If $F(x)$ is absolutely continuous then $f=F'$ exists a.e.~and $f\in L^1$ and $F(x)-F(a) = \int_a^x f(t)\,dt$ for all $x$.
\end{theorem}

\bigskip
\begin{itemize}
\item take-home message: if $f\in L^1$ then its indefinite integral $F(x) = \int_a^x f(t)\,dt$ is a nice ``very'' continuous function which is differentiable a.e.
\item \emph{and} you can recover $f$ from $F$ by differentiating, at least a.e.
\end{itemize}
\end{frame}


\begin{frame}{the devil's staircase}

\medskip \noindent \emph{question}: if $F:[a,b]\to\RR$ is continuous, and if its derivative $f=F'$ exists almost everywhere, and if $f\in L^1$, then is true that
    $$F(x) -F(a)= \int_a^x f(t)\,dt \quad \text{for all } x\in [a,b]?$$

\medskip \noindent \emph{answer}: \alert{no}

\bigskip
\noindent picture:

\vspace{50mm}
\end{frame}



% at start of sections 2,3,... put outline
\AtBeginSection[] {
    \begin{frame}{Outline}
    \setbeamertemplate{section in toc}[sections numbered]
    \tableofcontents[currentsection]%,hideallsubsections]
    \end{frame}
}


\section{integration by parts}

\begin{frame}{y}

\begin{itemize}
\item y
\end{itemize}
\end{frame}


\section{divergence theorem on $\RR^n$}

\begin{frame}{y}

\begin{itemize}
\item y
\end{itemize}
\end{frame}


\section{weak derivatives}

\begin{frame}{y}

\begin{itemize}
\item y
\end{itemize}
\end{frame}


\section{where this is going: Sobolev spaces to solve PDEs}

\begin{frame}{y}

\begin{itemize}
\item y
\end{itemize}
\end{frame}


\section{theory in weeks 9 \& 10}

\begin{frame}{theory in weeks 9 \& 10}

\begin{itemize}
\item part of this material is only from slides
\item \dots but part will be from Chapter 5 in Saxe, sections 5.1 and 5.2
\end{itemize}

\begin{block}{to define:}
\begin{itemize}
\item linear operator
\item bounded linear operator
\item operator norm
\item dual space
\item Sobolev space $H^1$
\end{itemize}
\end{block}
\end{frame}


\begin{frame}{theory in weeks 9 \& 10}

\begin{block}{to prove:}
\begin{itemize}
\item a linear operator is continuous if and only if it is bounded (Theorem 5.2)
\item the bounded linear operators $\cB(X,Y)$ form a normed vector space under operator norm (Theorem 5.3)
\item if $Y$ is complete then $\cB(X,Y)$ is complete (Theorem 5.4)
\item the dual space is a Banach space
\item dual vectors on a Hilbert space can be represented by elements of the Hilbert space (Riesz representation theorem)
\end{itemize}
\end{block}

\begin{itemize}
\item Assignments 5 and 6 are posted at \quad  \href{https://bueler.github.io/fa/index.html}{\texttt{bueler.github.io/fa}}
\end{itemize}
\end{frame}

\end{document}
